\documentclass[11pt]{article}

\usepackage[margin=1in]{geometry}
\usepackage[T1]{fontenc}
\usepackage[utf8]{inputenc}
\usepackage{lmodern}
\usepackage{microtype}
\usepackage{amsmath, amssymb}
\usepackage{booktabs}
\usepackage{hyperref}

\title{repx: Kernel-Level Process Attestation for Reproducible Build Provenance}
\author{Anonymous}
\date{\today}

\begin{document}
\maketitle

\begin{abstract}
Software supply-chain security has focused heavily on artifact metadata (for example, SBOMs), but many high-impact compromises occur in the \emph{build process} itself. This paper presents \texttt{repx}, a Linux prototype that generates cryptographic attestations of build-time process behavior by tracing kernel events with eBPF, canonicalizing them into system-independent operations, and hashing them into a Merkle-tree commitment. The main research value is methodological: \texttt{repx} shows how to turn low-level runtime traces into deterministic provenance summaries suitable for verification and divergence analysis. We describe the architecture, define the achieved guarantees, and identify open research problems (coverage, loss handling, and formal security properties). The contribution is not merely a tool implementation, but an executable design pattern for process-centric provenance in reproducible systems.
\end{abstract}

\section{Introduction}
Build reproducibility and software provenance are central topics in modern software assurance. Existing approaches often verify \emph{what} dependencies were declared or \emph{what} artifact was produced, but they are less effective at proving \emph{how} the artifact was produced at runtime. This gap is important because attacker behavior frequently manifests as extra file accesses, unauthorized process launches, or hidden outputs that may be invisible to source review and SBOM inspection.

\texttt{repx} addresses this gap by capturing process-level behavior at the kernel boundary and producing an attestation hash that can be checked later via re-execution. The project targets a core research question:
\begin{quote}
Can kernel-observed build behavior be canonicalized into a deterministic, verifiable process provenance commitment?
\end{quote}

The current implementation indicates that the answer is ``yes'' for a practical subset of build-relevant events on Linux.

\section{Problem Formulation}
Let a build run produce a time-ordered stream of observed runtime events
\[
E = (e_1, e_2, \dots, e_n),
\]
where each event is emitted from kernel tracepoints. The implementation maps this stream to a canonical set of distinct operations
\[
C = \operatorname{set}(\mathcal{N}(E)),
\]
where $\mathcal{N}$ normalizes system-dependent details and operation hashes are sorted and deduplicated. Counts, ordering, timestamps, and logical process indices are not part of the semantic commitment. The system then computes
\[
R_p = \mathcal{M}(C),
\]
where $\mathcal{M}$ is a Merkle-tree root over the sorted operation hashes. A top-level root additionally binds $R_p$ to the exact command, output-selection policy, and selected output hashes.

Two runs are equivalent under this model if and only if they produce the same top-level root. Because $C$ is a set rather than a sequence, diagnostics use set difference over leaf hashes; positional Merkle descent would produce cascading differences after an insertion.

\section{System Overview}
\texttt{repx} consists of three components:
\begin{enumerate}
  \item \textbf{eBPF collection layer} (kernel-space): captures selected tracepoints for open/close, mmap, exec, fork, and exit, scoped to a root process and descendants.
  \item \textbf{Userspace tracer and canonicalizer}: resolves successful opens through \texttt{/proc/<pid>/fd}, validates resolved targets against captured paths, retains nonblocking regular-file handles when possible, filters external watch events before observation, hashes file identity as events arrive, normalizes selected system-state accesses, and maps real PIDs to diagnostic logical indices.
  \item \textbf{Attestation layer}: hashes canonical operations and builds a Merkle tree; emits a JSON attestation containing the sorted leaves plus a root that binds the command and output policy.
\end{enumerate}

Verification re-runs the command, recomputes the root, and reports either equality (verified) or divergence points (diagnostics).

\section{What This Achieves Academically}
\subsection{A Process-Centric Provenance Primitive}
\texttt{repx} operationalizes ``Software Bill of Process'' style reasoning: provenance is derived from observed runtime behavior rather than declared intent. This shifts the measurement boundary from package metadata to kernel-observed execution, enabling investigation of attacks that manipulate build scripts, transient files, or process chains.

\subsection{A Canonicalization Strategy for Reproducibility}
The prototype contributes a concrete normalization pipeline:
\begin{itemize}
  \item event-arrival hashing of file inputs/outputs rather than end-of-run path hashing,
  \item unavailable identities that distinguish missing, unreadable, and non-regular files, bind stable paths, and normalize recognized session/compiler temporary names,
  \item output-rooted slicing that commits selected output hashes and content-resolved dependencies while omitting unavailable transient intermediates from the dependency walk,
  \item normalization of selected host-state reads (for example, \texttt{/proc}, \texttt{/sys}) via sentinel values,
  \item logical process indexing and root-process anchoring to reduce run-specific variance.
\end{itemize}
This is valuable as a reusable pattern for future provenance systems.

\subsection{Commit-and-Diff Verifiability}
The sorted-set Merkle commitment provides two academic benefits:
\begin{itemize}
  \item \textbf{compact verification}: one root hash for equality checks,
  \item \textbf{explainable failure}: added and missing operation hashes without positional cascades.
\end{itemize}
This combines cryptographic succinctness with operational diagnosability.

\subsection{Integrity-Aware Handling of Telemetry Loss}
The implementation tracks dropped ring-buffer events and fails closed by default:
\begin{itemize}
  \item trace mode refuses to emit an attestation unless the user explicitly supplies \texttt{--allow-dropped-events},
  \item verify mode fails closed on dropped events.
\end{itemize}
This explicit treatment of observability failure is a meaningful research property; many tracing systems silently degrade.

\section{Threat Model and Security Semantics}
\textbf{In-scope:} behavioral changes that alter traced process activity (extra writes, modified tool binaries, changed command paths, altered build scripts, new process chains, etc.).

\textbf{Out-of-scope / partial:}
\begin{itemize}
  \item attacks entirely outside traced event coverage,
  \item kernel compromise (the monitor and measured environment are not independent),
  \item races where \texttt{/proc}-based best-effort enrichments are unavailable,
  \item semantics not yet encoded in canonicalization (for example, complete filesystem namespace resolution),
  \item replacement of the attestation file itself: the JSON document is not yet signed or wrapped in DSSE.
\end{itemize}

Therefore, the attestation is best interpreted as a \emph{deterministic commitment over observed covered behavior}, not a full proof of non-malicious execution.

\subsection{Explicit Evasion and Coverage Table}
The prototype's collection boundary is security-relevant. An adversary controlling the build can deliberately select an uncovered mechanism.

\begin{center}
\begin{tabular}{p{0.28\linewidth}p{0.62\linewidth}}
\toprule
Mechanism & Current status \\
\midrule
\texttt{openat}, \texttt{close}, file-backed \texttt{mmap} & Covered; \texttt{O\_RDWR} is modeled as both read and write. Writable private mappings are reads, while only writable \texttt{MAP\_SHARED} mappings are writes. \\
\texttt{open}, \texttt{openat2} & Not traced. \\
\texttt{rename*}, \texttt{unlink*} & Successful operations are traced. Descriptor paths and output-slice writer aliases follow final rename destinations; deletion remains a distinct operation. \\
\texttt{dup*}, inherited descriptors & Descriptor attribution is incomplete. \\
\texttt{sendfile}, \texttt{copy\_file\_range}, \texttt{io\_uring} & Not traced. \\
Network I/O & Not traced; remote scripts or data are outside the commitment unless they also appear through a covered file access. \\
Watch-mode relative reads & May be missed because the eBPF prefix check sees the raw relative pathname before userspace resolution. \\
\bottomrule
\end{tabular}
\end{center}

This table is part of the security claim, not merely an implementation backlog.

\section{Prototype Evidence}
The repository includes integration scenarios demonstrating:
\begin{enumerate}
  \item successful trace/verify consistency on deterministic C build workflows,
  \item detection of process-level divergence in a script-tampering ``attack demo'' where hidden output is introduced.
\end{enumerate}
These are prototype-level case studies, not a statistically powered benchmark. Their research value is proof-of-feasibility for the process-attestation model.

\section{Limitations (Important for Correct Interpretation)}
\begin{itemize}
  \item \textbf{Platform scope:} Linux + eBPF, requiring elevated privileges/capabilities.
  \item \textbf{Coverage scope:} only selected syscalls/tracepoints are modeled.
  \item \textbf{Set semantics:} repeated operations and temporal order are intentionally discarded for reproducibility.
  \item \textbf{Unsigned envelope:} the root binds command and policy fields, but authenticity still depends on trusted distribution of either the JSON file or an out-of-band root supplied with \texttt{--expected-root}.
  \item \textbf{Canonicalization assumptions:} some path and process metadata enrichment is best-effort.
  \item \textbf{No formal proof yet:} the current guarantees are empirical and implementation-driven.
  \item \textbf{No large-scale evaluation yet:} overhead and false-positive/negative behavior across diverse ecosystems remain open.
\end{itemize}

\section{Research Opportunities Enabled by repx}
\begin{enumerate}
  \item \textbf{Formal semantics of canonicalization:} prove invariants and equivalence classes for reproducible runs.
  \item \textbf{Coverage optimization:} identify minimal tracepoint sets that maximize attack detectability.
  \item \textbf{Cross-host reproducibility studies:} quantify which host differences must be normalized.
  \item \textbf{Compositional attestations:} combine process commitments with SBOM/in-toto/SLSA predicates.
  \item \textbf{Adversarial evaluation:} systematically test evasion strategies against event coverage and normalization rules.
\end{enumerate}

\section{Conclusion}
\texttt{repx} demonstrates a practical and academically relevant pattern: selected kernel-level execution observations can be transformed into deterministic, hash-committed process provenance. The strongest contribution is conceptual and methodological: a concrete pathway from runtime telemetry to reproducible attestations with set-based divergence diagnostics. The prototype does not prove that all build behavior was observed and does not yet authenticate its JSON envelope, but it provides a defensible foundation for research on process-centric supply-chain integrity and reproducible build assurance.

\begin{thebibliography}{9}

\bibitem{merkle}
R. C. Merkle.
\newblock A Digital Signature Based on a Conventional Encryption Function.
\newblock \emph{Advances in Cryptology---CRYPTO '87}, 1988.

\bibitem{slsa}
SLSA Authors.
\newblock Supply-chain Levels for Software Artifacts (SLSA).
\newblock \url{https://slsa.dev/}

\bibitem{intoto}
in-toto Authors.
\newblock in-toto: A framework to secure the integrity of software supply chains.
\newblock \url{https://in-toto.io/}

\bibitem{ebpf}
eBPF Foundation.
\newblock What is eBPF?
\newblock \url{https://ebpf.io/what-is-ebpf/}

\end{thebibliography}

\end{document}
